\chapter{Исследовательский раздел}

% TODO красиво оформить список источников
% TODO Отладить код так, чтобы было что показать на защите
% TODO Презентацию сделать
%% TODO Добавить датасет с людьми в часть с кодом

Цель данного исследования --- объективная оценка качества реконструкции глубины посредством разработанного гибридного метода по сравнению с современными подходами, с особым вниманием к устойчивости методов при работе с деградировавшими изображениями. Исследование направлено на определение сильных и слабых сторон разработанного метода, а также его применимости к архивным фотографиям с различными характеристиками.

Основные задачи исследования:
\begin{itemize}
    \item проанализировать влияние различных типов деградации изображений на качество реконструкции;
    \item сравнить точность восстановления глубины разработанного метода с современными подходами на эталонных наборах данных;
    \item определить наиболее подходящие области применения каждого метода.
\end{itemize}

\section{Описание входных данных}

Ввиду отсутствия эталонных данных о глубине для архивных фотографий,
разработан методологический подход с использованием современных датасетов с верифицированными картами глубины и алгоритмов искусственного старения,
позволяющий количественно оценивать результаты реконструкции в контролируемых условиях.

Для проведения исследования были выбраны следующие датасеты, обеспечивающих разнообразие тестовых сцен:

\begin{enumerate}
    \item \textbf{NYU Depth v2} --- 1449 изображений внутренних помещений с глубиной, полученной с помощью сенсора Kinect~\cite{silberman2012indoor}.
    \item \textbf{4D Light Field Benchmark} --- 24 синтетические сцены с глубиной, полученной с помощью рендеринга 3D-моделей~\cite{honauer2017dataset}.
    \item \textbf{BIWI Kinect Head Pose} --- набор из более 15000 изображений 20 человек, полученных с помощью сенсора Kinect~\cite{fanelli_IJCV}.
\end{enumerate}

Из каждого датасета было случайным образом отобрано по 20 изображений,
обеспечивая статистическую значимость результатов при сохранении выполнимости тестирования в разумные сроки.

Для имитации состаривания изображений был разработан программный модуль, реализующий следующие эффекты с параметрически задаваемой интенсивностью (от 0 до 1):

\begin{itemize}
     \item \textbf{сепия} --- преобразование RGB в тонированное монохромное изображение,
     \item \textbf{виньетирование} --- затемнение краев изображения,
     \item \textbf{шум и зернистость} --- добавление шума, имитирующего фотопленку,
     \item \textbf{выцветание} --- снижение контраста и насыщенности,
     \item \textbf{случайные пятна и царапины} --- добавление локальных дефектов.
\end{itemize}

\section{Метрики оценки качества}

% TODO надо ли описание элементоов формулы оформлять в виде списка или и так сойдет?

Для объективной оценки качества восстановления глубины используются метрики, которые позволяют количественно оценить,
насколько точно модель предсказывает карту глубины, сравнивая предсказанные значения с истинными.

Для корректного сравнения методов, формирующих результаты в разных масштабах, в работе используется предварительная нормализация карт глубин, которая выполняется по формуле:
\begin{equation}
    d_{\text{norm}} = \frac{d - d_{\text{min}}}{d_{\text{max}} - d_{\text{min}}},
\end{equation}
$d_{\text{norm}}$ --- нормализованное значение глубины, 
$d$ --- исходное значение глубины, 
$d_{\text{min}}$ --- минимальное значение глубины в карте глубин, 
$d_{\text{max}}$ --- максимальное значение глубины в карте глубин. 

Для определения среднего отклонения результата используется метрика $RMSE$ --- среднеквадратичное отклонение результата
(от англ. Root Mean Square Error), которая позволяет оценить,
насколько предсказанные значения глубины отличны от истинных значений в среднем по всем пикселям изображения~\cite{zhang2025review, istomin2024, shestopalov2024}.
$RMSE$ вычисляется согласно формуле:

\begin{equation}
    \text{RMSE} = \sqrt{\frac{1}{n} \sum_{i=1}^{n} (d_{i}^{\text{pred}} - d_{i}^{\text{gt}})^2},
\end{equation}
 где $n$ -- общее количество пикселей изображения,
$d_{i}^{\text{pred}}$ -- предсказанное значение глубины для $i$-го пикселя,
$d_{i}^{\text{gt}}$ -- истинное (эталонное) значение глубины для $i$-го пикселя.

С учетом нормализации, минимальное значение RMSE --- 0, максимальное --- 1.
Чем ближе значение метрики к 0, тем точнее предсказание.

С целью определения количества предсказанных значений глубины, которые находятся в пределах допустимого отклонения используется
пороговая относительная точность ($\delta_{\lambda}$)~\cite{zhang2025review, istomin2024, shestopalov2024}.
Эта метрика измеряет долю пикселей, для которых предсказанные и истинные значения глубины удовлетворяют следующему условию:

\begin{equation} 
    \max\left(\frac{d_{i}^{\text{pred}}}{d_{i}^{\text{gt}}}, \frac{d_{i}^{\text{gt}}}{d_{i}^{\text{pred}}}\right) = \delta < t_{\lambda} 
\end{equation}

где $d_{i}^{\text{gt}}$ --- истинное значение глубины для $i$-го пикселя,
$d_{i}^{\text{pred}}$ --- предсказанное значение глубины для $i$-го пикселя,
$\delta$ --- значение, показывающее отклонение предсказанного значения от истинного,
$t_{\lambda}$ --- пороговое значение, при котором отклонение считается допустимым.

В работе используются пороги $\lambda=1,2,3$, что соответствует отклонению в 25\%, 56.25\% и 95.3\% соответственно.

С учетом нормализации значения метрики $\delta_{\lambda}$ лежат в диапазоне от 0 до 1.
Чем ближе значение метрики к 1, тем большая доля пикселей имеет точное предсказание глубины в пределах соответствующего порога.

\section{Исследование устойчивости метода к деградации изображений}

Исследуется эффективность гибридного метода в зависимости от интенсивности различных эффектов старения,
применённых к тестовым изображениям изолированно.

Для каждого эффекта (сепия, виньетирование, шум, выцветание, пятна) применялись различные уровни интенсивности от 0.0 до 1.0 с шагом 0.2.

Гипотеза заключается в том, что различные эффекты старения будут по-разному влиять на точность реконструкции глубины:
эффекты, изменяющие только цветовую гамму изображения (сепия, выцветание), окажут меньшее влияние на точность,
чем эффекты, искажающие структурную информацию изображения (шум, пятна, виньетирование).

На рисунке~\ref{fig:effects-rmse} представлены графики зависимости среднеквадратичной ошибки от интенсивности применения различных эффектов старения к тестовым изображениям.

\begin{figure}
\centering
    \begin{tikzpicture}
        \begin{axis}[
            width=0.9\textwidth,
            height=6cm,
            xlabel={Интенсивность эффекта},
            ylabel={Среднеквадратичная ошибка (RMSE)},
            xmin=-0.05, xmax=1.05,
            ymin=0.21, ymax=0.33,
            xtick={0,0.2,0.4,0.6,0.8,1.0},
            grid=both,
            grid style={line width=.1pt, draw=gray!10},
            major grid style={line width=.2pt,draw=gray!50},
            every axis plot/.append style={smooth},
            legend pos=north west,
            legend style={font=\small},
            legend columns=2]

            % Сепия
            \addplot[black, solid, smooth] coordinates {
                (0.0, 0.22673959517265008)
                (0.2, 0.2224120524354949)
                (0.4, 0.22380517749733062)
                (0.6, 0.22040347321392373)
                (0.8, 0.22158114785240235)
                (1.0, 0.22479423061609816)
            };

            % Виньетирование
            \addplot[black, dashed, smooth] coordinates {
                (0.0, 0.22673959517265008)
                (0.2, 0.256956901702622)
                (0.4, 0.2672339682435893)
                (0.6, 0.27263466877343556)
                (0.8, 0.2877002885964059)
                (1.0, 0.2835513838997222)
            };

            % Шум
            \addplot[black, dotted, smooth, line width=1.25pt] coordinates {
                (0.0, 0.22673959517265008)
                (0.2, 0.22550388966890234)
                (0.4, 0.2474848386682945)
                (0.6, 0.312124558133447)
                (0.8, 0.32290318666871687)
                (1.0, 0.2973997438450658)
            };

            % Выцветание
            \addplot[black, dashdotted, smooth] coordinates {
                (0.0, 0.22673959517265008)
                (0.2, 0.22192686388303162)
                (0.4, 0.21995908234551922)
                (0.6, 0.22126814934419525)
                (0.8, 0.23044226142223376)
                (1.0, 0.22065211351964661)
            };

            % Пятна
            \addplot[black, loosely dashed, smooth] coordinates {
                (0.0, 0.22673959517265008)
                (0.2, 0.24074089383533015)
                (0.4, 0.23506201341607524)
                (0.6, 0.2461171888472462)
                (0.8, 0.24989895793232644)
                (1.0, 0.2392362046372416)
            };

            \legend{Сепия, Виньетирование, Шум, Выцветание, Пятна}
        \end{axis}
    \end{tikzpicture}
    \caption{Влияние эффектов старения на среднеквадратичную ошибку (RMSE)}
    \label{fig:effects-rmse}
\end{figure}

Анализ полученных данных подтверждает выдвинутую гипотезу о том, что эффекты, изменяющие только цветовую гамму изображения (сепия, выцветание),
оказывают меньшее влияние на точность реконструкции, чем эффекты, искажающие структурную информацию (шум, виньетирование, пятна).

\begin{enumerate}
    \item \textbf{Сепия и выцветание} оказывают минимальное влияние на точность метода:
    при максимальной интенсивности эффектов для сепии отклонение RMSE составляет 0.0006 (0.26\%), для выцветания — 0.004 (1.6\%).
    Это объясняется тем, что преобразования цветовой гаммы не изменяют структурную информацию изображения, а алгоритмы восстановления глубины опираются преимущественно на структурные, а не цветовые характеристики.

    \item \textbf{Виньетирование} демонстрирует значительно более выраженное влияние:
    при интенсивности 0.8 наблюдается максимальное увеличение RMSE на 0.061 (26.8\%).
    Затемнение по краям изображения существенно влияет на восприятие глубины, особенно в периферийных областях, что подтверждает гипотезу о важности равномерного освещения для структурного анализа изображения.

    \item \textbf{Шум} показывает наиболее сильное влияние на точность реконструкции:
    при интенсивности 0.8 RMSE возрастает на 0.096 (42.4\%).
    Шум разрушает текстуры и размывает границы объектов, что критически затрудняет анализ формы из затенения и работу нейросетевых алгоритмов.

    \item \textbf{Пятна и царапины} демонстрируют умеренное, но устойчивое влияние на точность:
    максимальное увеличение RMSE на 0.023 (10.2\%) наблюдается при интенсивности 0.8.
    Локальные искажения структуры изображения влияют на алгоритм менее критично, чем равномерный шум, но более существенно, чем изменения цветовой гаммы.
\end{enumerate}

Результаты позволяют определить наиболее эффективные области применения разработанного метода:
реконструкция трёхмерной информации из фотографий с преимущественно цветовой деградацией (сепия, выцветание) и умеренными локальными дефектами.
В то же время для изображений с сильным шумом или виньетированием необходима более тщательная предварительная обработка.

\section{Сравнительный анализ методов}

В данном исследовании сравнивалась три метода восстановления глубины из монокулярных изображений:
\begin{itemize}
    \item \textbf{восстановление глубины по закраске изображения (SFS)} --- метод, основанный на анализе формы из затенения (Shape From Shading), использующий информацию о характере освещения для восстановления формы объектов;
    \item \textbf{MiDaS} --- метод восстановления глубины на основе глубоких нейронных сетей, обученный на разнородных данных для повышения обобщающей способности;
    \item \textbf{гибридный метод} --- разработанный в рамках исследования подход, адаптивно интегрирующий нейросетевой метод MiDaS и классический метод Shape From Shading.
\end{itemize}

С целью комплексного исследования влияния деградации изображений на качество восстановления глубины
применялась методика контролируемого искусственного старения.
Эффекты деградации моделировались с различными интенсивностями: для виньетирования и пятен — в диапазоне от 0 до 0.6, а для сепии и выцветания — от 0 до 1.0, что основано на результатах предыдущего исследования.

Гипотеза заключается в том, что гибридный метод будет объединять преимущества обоих подходов:
стабильность MiDaS при работе с деградированными изображениями и высокую детализацию SFS на качественных изображениях. 
Ожидается, что при увеличении степени деградации изображений гибридный метод будет демонстрировать большую устойчивость, 
чем чистый SFS, сохраняя при этом лучшую точность восстановления геометрических деталей по сравнению с чистым MiDaS.

В таблице~\ref{tab:comparison-by-degradation} представлены результаты сравнения трех методов по метрикам RMSE и $\delta_{\lambda}$ при различных уровнях деградации изображений.

\begin{table}[H]
\centering
\caption{Сравнение методов восстановления глубины при различных уровнях деградации изображений}
\label{tab:comparison-by-degradation}
\begin{tabular}{|l|c|c|c|c|c|}
\hline
\textbf{Метод} & \textbf{Уровень} & \textbf{RMSE} & \textbf{$\delta_1$} & \textbf{$\delta_2$} & \textbf{$\delta_3$} \\
\hline
\multirow{Гибридный} & 0 & 0.2854 & 0.6376 & 0.8739 & 0.9760 \\
 & 1 & 0.2954 & 0.6380 & 0.8666 & 0.9677 \\
 & 2 & 0.2931 & 0.6849 & 0.8863 & 0.9667 \\
 & 3 & 0.2918 & 0.6415 & 0.8638 & 0.9565 \\
\hline
\multirow{MiDaS} & 0 & 0.2914 & 0.6425 & 0.8597 & 0.9542 \\
 & 1 & 0.2947 & 0.6399 & 0.8607 & 0.9542 \\
 & 2 & 0.2983 & 0.6715 & 0.8735 & 0.9522 \\
 & 3 & 0.2979 & 0.6438 & 0.8646 & 0.9427 \\
\hline
\multirow{SFS} & 0 & 0.2867 & 0.6768 & 0.8866 & 0.9790 \\
 & 1 & 0.3832 & 0.3664 & 0.5749 & 0.7588 \\
 & 2 & 0.3959 & 0.3749 & 0.6144 & 0.7667 \\
 & 3 & 0.4170 & 0.4133 & 0.6090 & 0.7622 \\
\hline
\end{tabular}
\end{table}

На рисунках~\ref{fig:comparison-rmse}, ~\ref{fig:comparison-by-degradation} представлены графики, визуализирующие результаты исследования.

\begin{figure}[H]
\centering
    \begin{tikzpicture}
        \begin{axis}[
            width=\textwidth,
            height=6cm,
            xlabel={Уровень деградации},
            ylabel={Среднеквадратичная ошибка (RMSE)},
            xmin=-0.1, xmax=3.1,
            ymin=0.25, ymax=0.6,
            xtick={0,1,2,3},
            ytick={0.3,0.35,0.4,0.45,0.5,0.55,0.6},
            grid=both,
            grid style={line width=.1pt, draw=gray!10},
            major grid style={line width=.2pt,draw=gray!50},
            legend pos=north east,
            legend style={font=\small}]

            % Гибридный метод
            \addplot[black, mark=*, smooth] coordinates {
                (0, 0.28541211159480284)
                (1, 0.295440172140253)
                (2, 0.2931316731338902)
                (3, 0.2918331874608755)
            };

            % MiDaS
            \addplot[black, dashed, mark=square*, smooth] coordinates {
                (0, 0.29138717467995495)
                (1, 0.2947147948982483)
                (2, 0.29826661689549494)
                (3, 0.2979285750148052)
            };

            % SFS (Segmentation)
            \addplot[black, dotted, mark=triangle*, smooth] coordinates {
                (0, 0.2867484161955488)
                (1, 0.3832440906946321)
                (2, 0.39591868578935213)
                (3, 0.4169904701291959)
            };

            \legend{Гибридный метод, MiDaS, SFS}
        \end{axis}
    \end{tikzpicture}
    \caption{Сравнение методов по метрике RMSE при различных уровнях деградации}
    \label{fig:comparison-rmse}
\end{figure}

\begin{figure}[H]
\centering
    \begin{tikzpicture}
        \begin{axis}[
            width=\textwidth,
            height=8cm,
            xlabel={Уровень деградации},
            ylabel={Значение метрик $\delta_{\lambda}$},
            xmin=-0.5, xmax=3.5,
            ymin=0.35, ymax=1,
            xtick={0,1,2,3},
            ytick={0.4,0.5,0.6,0.7,0.8,0.9, 1},
            grid=both,
            grid style={line width=.1pt, draw=gray!10},
            major grid style={line width=.2pt,draw=gray!50},
            legend style={font=\small,at={(0.5,-0.25)},anchor=north,legend columns=3},
            legend cell align={left}]

            % Гибридный метод - Delta1
            \addplot[black, mark=*, smooth] coordinates {
                (0, 0.6376040039062501)
                (1, 0.6380177408854165)
                (2, 0.6849495442708332)
                (3, 0.6414757486979165)
            };

            % MiDaS - Delta1
            \addplot[black, dashed, mark=*, smooth] coordinates {
                (0, 0.6425073242187499)
                (1, 0.6399417317708334)
                (2, 0.6715174153645833)
                (3, 0.6438002929687499)
            };

            % SFS - Delta1
            \addplot[black, dotted, mark=*, smooth] coordinates {
                (0, 0.6767942708333333)
                (1, 0.36636621093749996)
                (2, 0.3749396158854167)
                (3, 0.41326790364583335)
            };

            % Гибридный метод - Delta2
            \addplot[black, mark=square*, smooth] coordinates {
                (0, 0.8739314778645833)
                (1, 0.8665922851562499)
                (2, 0.8863361002604167)
                (3, 0.8638222656250001)
            };

            % MiDaS - Delta2
            \addplot[black, dashed, mark=square*, smooth] coordinates {
                (0, 0.8597228190104166)
                (1, 0.8607184244791665)
                (2, 0.8735188802083333)
                (3, 0.8645610351562502)
            };

            % SFS - Delta2
            \addplot[black, dotted, mark=square*, smooth] coordinates {
                (0, 0.8865755208333334)
                (1, 0.57492626953125)
                (2, 0.6144095052083333)
                (3, 0.6090374348958332)
            };

            % Гибридный метод - Delta3
            \addplot[black, mark=triangle*, smooth] coordinates {
                (0, 0.9759777018229167)
                (1, 0.9676751302083334)
                (2, 0.9666614583333334)
                (3, 0.9564737955729166)
            };

            % MiDaS - Delta3
            \addplot[black, dashed, mark=triangle*, smooth] coordinates {
                (0, 0.9542218424479169)
                (1, 0.9541813151041666)
                (2, 0.9522010091145832)
                (3, 0.9427171223958334)
            };

            % SFS - Delta3
            \addplot[black, dotted, mark=triangle*, smooth] coordinates {
                (0, 0.9790131835937501)
                (1, 0.7588245442708333)
                (2, 0.76665478515625)
                (3, 0.7621984049479166)
            };

            \legend{$\delta_1$ Гибридный, $\delta_1$ MiDaS, $\delta_1$ SFS,
                    $\delta_2$ Гибридный, $\delta_2$ MiDaS, $\delta_2$ SFS,
                    $\delta_3$ Гибридный, $\delta_3$ MiDaS, $\delta_3$ SFS}
        \end{axis}
    \end{tikzpicture}
    \caption{Сравнение методов по метрикам $\delta_{\lambda}$ при различных уровнях деградации}
    \label{fig:comparison-by-degradation}
\end{figure}

На качественных изображениях (уровень деградации 0) все три метода показывают сопоставимые результаты по метрике RMSE: 
гибридный метод (0.2854), SFS (0.2867) и MiDaS (0.2914).
При этом SFS демонстрирует наилучшие показатели по метрикам $\delta_1$, $\delta_2$ и $\delta_3$.

При увеличении уровня деградации (уровни 1-3) наблюдается следующая динамика:
\begin{itemize}
    \item метод SFS критически теряет в точности: RMSE возрастает на 45.4\%, а метрики $\delta_1$, $\delta_2$ и $\delta_3$ падают на 39-41\%;
    \item MiDaS демонстрирует стабильные показатели: RMSE практически не изменяется (0.2947-0.2983), как и метрики $\delta_\lambda$;
    \item гибридный метод сохраняет устойчивость на уровне MiDaS, но с лучшими абсолютными показателями:
RMSE в диапазоне 0.2918-0.2954 и метрики $\delta_2$ и $\delta_3$ в пределах 0.8638-0.8863 и 0.9565-0.9677 соответственно.
\end{itemize}

Результаты сравнительного анализа трех методов восстановления глубины подтверждают выдвинутую гипотезу: по сравнению с базовыми методами
гибридный демонстрирует незначительное повышение точности (около 1-2\%), которое стабильно наблюдается на всех уровнях деградации изображений.


Полученные результаты позволяют определить оптимальные области применения каждого из методов.
\begin{itemize} \item \textbf{SFS} наиболее эффективен при работе с качественными изображениями без деградации, где обеспечивает высокую детализацию поверхностей. Рекомендуется для применения в контролируемых условиях, например, в промышленных системах контроля качества или при работе с современными цифровыми изображениями высокого разрешения;
\item \textbf{MiDaS} демонстрирует универсальность и устойчивость к различным типам искажений, что делает его подходящим для обработки разнородных коллекций изображений. Оптимален для задач, где важна стабильность результатов, но высокая детализация не является критичной;
\item \textbf{Гибридный метод} обеспечивает оптимальный баланс между устойчивостью к деградации и детализацией при восстановлении трёхмерной информации из архивных фотографий. Метод повышает точность реконструкции без необходимости переобучения нейросетевых моделей, что существенно снижает вычислительные затраты.
\end{itemize}

\section*{Вывод}

% TODO переписать

В исследовательском разделе проведено исследование с использованием стандартных датасетов,
алгоритмов искусственного старения изображений и объективных метрик оценки
для количественной оценки эффективности разработанного гибридного метода восстановления глубины.

Проведено исследование устойчивости гибридного метода к различным эффектам старения изображений,
которое подтвердило гипотезу о дифференцированном влиянии различных типов деградации на точность реконструкции глубины:
    \begin{itemize}
        \item эффекты, изменяющие структурную информацию (шум и виньетирование), оказывают наиболее сильное негативное влияние (26.8-42.4\% увеличения RMSE);
        \item эффекты, затрагивающие только цветовую гамму (сепия и выцветание), влияют минимально (0.26-1.6\% увеличения RMSE);
        \item локальные искажения (пятна и царапины) оказывают умеренное влияние (около 10.2\% увеличения RMSE).
    \end{itemize}

Результаты позволяют определить наиболее эффективные области применения разработанного метода:
реконструкция трёхмерной информации из фотографий с преимущественно цветовой деградацией (сепия, выцветание) и умеренными локальными дефектами.
В то же время для изображений с сильным шумом или виньетированием необходима более тщательная предварительная обработка.

Выполнен сравнительный анализ разработанного гибридного метода с базовыми подходами (MiDaS и SFS) при различных уровнях деградации изображений:
    \begin{itemize}
        \item на качественных изображениях все три метода демонстрируют сопоставимые результаты по метрике RMSE;
        \item при увеличении уровня деградации метод SFS критически теряет в точности (RMSE возрастает на 45.4\%);
        \item MiDaS сохраняет стабильность на всех уровнях деградации;
        \item гибридный метод демонстрирует устойчивость на уровне MiDaS с улучшенными на 1-2\% абсолютными показателями точности.
    \end{itemize}

Полученные результаты позволяют определить оптимальные области применения каждого из методов.
\begin{itemize} \item \textbf{SFS} наиболее эффективен при работе с качественными изображениями без деградации, где обеспечивает высокую детализацию поверхностей. Рекомендуется для применения в контролируемых условиях, например, в промышленных системах контроля качества или при работе с современными цифровыми изображениями высокого разрешения;
\item \textbf{MiDaS} демонстрирует универсальность и устойчивость к различным типам искажений, что делает его подходящим для обработки разнородных коллекций изображений. Оптимален для задач, где важна стабильность результатов, но высокая детализация не является критичной;
\item \textbf{Гибридный метод} обеспечивает оптимальный баланс между устойчивостью к деградации и детализацией при восстановлении трёхмерной информации из архивных фотографий. Метод повышает точность реконструкции без необходимости переобучения нейросетевых моделей, что существенно снижает вычислительные затраты.
\end{itemize}

Результаты исследования подтверждают эффективность разработанного гибридного метода,
который успешно сочетает в себе устойчивость нейросетевого подхода к деградации изображений с детализирующими возможностями метода Shape from Shading.

Полученные данные позволяют определить оптимальные области применения метода:
восстановление трёхмерной информации из архивных фотографий для задач цифровой реставрации и трёхмерной реконструкции исторических объектов,
когда доступны только архивные фотоматериалы различного качества, но с минимальным уровнем шумовых искажений.
Метод позволяет получить более точные и детализированные карты глубины по сравнению с существующими подходами.


